\section{Ten unrestricted integers}
\label{sec:ten}

Supply exactly ten unrestricted integer coordinates
\[
 (f,g,h,k,l,w,x,y,z,m)\in\Z^{10}.
\]
Use the source of Section~\ref{sec:channels}, with
\begin{gather*}
 b=1+3f,\qquad B=a_Sb^2,\qquad \mathcal C=bB^{108},\qquad L_c=B^{2241},\\
 X=L_c^2+(S_c+T_c)L_c+L_c-1-S_c,\\
 Y=a_S^{15}\,2^{\kappa-21}B^{2356}b.
\end{gather*}
The monomial cap is \(\mathcal C=a_S^{108}b^{217}\), as required
by the inherited Sun bridge.  The source output \(Y\) is a fixed
polynomial, not an exponential with a variable exponent.

\subsection{The arithmetic source contract}

Before recovering a binary base, the reverse bridge has \(f\ne0\),
\(b\) an integer square, and \(0\le g<2\mathcal C\).  Since
\(b\equiv1\pmod3\), these hypotheses imply \(b\ge4\) and \(f>0\).
The preceding source proof applies and gives
\(S_c,T_c\ge0\) and \(S_c+T_c<L_c\), including the stronger summed
bound~\eqref{eq:source-sum-range}.  Hence
\[
 X\ge L_c^2>3b,\qquad Y>\max(b,256).
\]
The exponent \(\kappa\) is even, so \(B\equiv N_0\equiv L_c\equiv1\pmod3\).
We have \(M_c\equiv0\), \(A_T\equiv22\), and
\(T_c\equiv-22\pmod3\).  Consequently
\[
 X\equiv1+T_c\equiv0\pmod3.
\]
These range and congruence facts do not assume \(b\) is a power of two.

After binary-base recovery, write \(b=2^s\), \(B=2^v\).  For
nonnegative \(S,T\) with \(S+T<L_c\), the three disjoint blocks in
\(L_c^2+(S+T)L_c+L_c-1-S\) give
\[
 \pop(X)=1+\log_2L_c+\pop(T_c)
       -\bigl(\pop(S_c)+\pop(T_c)-\pop(S_c+T_c)\bigr).
\]
The subtracted carry deficit is nonnegative, and vanishes exactly when
\(S_c\mathbin{\&}T_c=0\).  The exact mask population is
\(\pop(T_c)=131v-31s-22\), so
\[
 1+\log_2L_c+\pop(T_c)=2372v-31s-21=\log_2Y.
\]
Legendre's formula therefore proves
\[
 Y\mid\binom{2X}{X}\quad\Longleftrightarrow\quad
 S_c\mathbin{\&}T_c=0.
\]
Conversely a parent zero can be encoded with arbitrarily large
\(b=4^s\), putting \(f=(b-1)/3\).  It has \(b\le g<\mathcal C\)
and the required divisibility.  This is the full source contract; the
Sun bridge does not require \(Y\) to be a square or equal to \(L_c^2\).

\subsection{Literal final polynomial}

The following finalizer is inherited from the inspected package-31 and
wave-7 construction; only its source polynomials have changed.  Define
\begin{gather*}
 L=lY,\quad U=2LX,\quad V=4gwY,\quad W_b=bw,\\
 A=U(V+1),\quad B_*=2X+1,\quad C=B_*+(A-2)h,\\
 P=U^2V,\quad K=X+1+k(P-2),\quad S_*=2A-5,\\
 T_*=3W_bC-2(W_b^2-1),\quad R_2=(P^2-4)K^2+4,\\
 \Delta_*=8\mathcal C^3gK^2-g^2\{32(C-KL)^2\mathcal C^3+g^2K^2\},\\
 O=f^2l^2\Delta_*,\quad\xi=zx,\quad D=(A^2-4)C^2+4,\\
 E=C^2D\xi,\quad F=4(A^2-4)E^2+1,\\
 G_*=2(A+2)(A-2)^2E^2-1,\quad\Omega_*=16E,\\
 H_*=C+(B_*-C)F+\Omega_*Fy,\quad I_*=(G_*^2-1)H_*^2+1.
\end{gather*}
Put
\begin{align*}
 r&=T_*+S_*y,&\Sigma&=\xi S_*^5,\\
 \Theta_*&=(r^3-rS_*^2-S_*^3)(r^2+23S_*^2),\\
 a_0&=b\{1+(3O-4)b^2z^2\},&a_1&=R_2,&a_2&=DFI_*,\\
 w_0&=1,&w_1&=a_0,&w_2&=a_0(P^2+K^2+3),\\
 Z_*&=\Sigma m-\Theta_*,&A_i&=\Sigma^2w_i^2a_i\quad(0\le i\le2),\\
 V_*&=Z_*^2+A_2-A_0-A_1,&
 W_*&=V_*^2+4Z_*^2A_2-4A_0A_1.
\end{align*}
The single equation is
\begin{equation}
 \boxed{F_{10}=W_*^2-16V_*^2Z_*^2A_2=0.}
 \label{eq:ten-final}
\end{equation}
All powers in these formulas are fixed nonnegative integers.  In
particular, no square root, factorial, binomial coefficient, bit operation,
or divisibility predicate occurs in the actual polynomial.

\subsection{Why the inherited finalizer still applies}

Sun's arithmetic bridge, specialized as in the preserved proof, supplies
both directions for any source satisfying the contract above.  The reverse
direction requires \(b\) square, \(O>0\), \(\xi\ne0\),
\(R_2\) and \(DFI_*\) square, and \(S_*\mid T_*\).  It recovers
\(b\) as a power of two, \(0<g<2\mathcal C\), and the central-binomial
divisibility.  The forward direction supplies these conditions, together
with the residue interpolation needed for the finalizer.  These are the
same arithmetic bridge hypotheses as in the inspected construction.

For clarity, the finalizer's main logical steps are recorded here.
The polynomial \((v^3-v-1)(v^2+23)\) has a root modulo every positive
integer and no rational root.  Prime valuations in its homogeneous
factors show that \(\Sigma\mid\Theta_*\) implies \(S_*\mid T_*\).
The parity of \(S_*\) makes it nonzero, and the absence of rational
roots makes \(\Theta_*\ne0\).  If \(\Sigma=0\), then
\(F_{10}=\Theta_*^8\ne0\), excluding that branch.

For \(z\ne0\), the factors \(b\) and
\(1+(3O-4)b^2z^2\) are coprime.  If their product is a square,
they are signed squares; \(b\equiv1\pmod3\) excludes a negative
square and the zero-factor exception.  The positivity Pell lemma then
gives \(O>0\).  Conversely, the forward bridge has \(O\ge4\);
a suitable positive Pell unit makes \(a_0\) a square.  Residue
interpolation supplies a multiple \(\xi=zx\) and the required
class of \(y\), making \(a_2\) a square and
\(\Sigma\mid\Theta_*\).

The structural weights satisfy
\[
 |w_1|>\sqrt{|a_0|},\qquad
 |w_2|>\sqrt{|a_0|}+|w_1|\sqrt{|a_1|}.
\]
Here \(|a_0|>1\) whenever \(z\ne0\), and
\(\sqrt{|a_1|}\le P^2+K^2+2\).  In the multiquadratic field,
any nontrivial sign-changing automorphism of a weighted root sum has
a largest changed index whose magnitude exceeds the sum of all smaller
changed magnitudes.  Thus a rational weighted sum forces every radicand
to be an integer square.  The exact identity
\[
 F_{10}=\prod_{\epsilon\in\{\pm1\}^3}
 \left(Z_*+\sum_{i=0}^2\epsilon_i\sqrt{A_i}\right)
\]
now proves the finalizer equivalence.  This argument also handles
negative radicands using complex absolute values, and zero radicands
separately.  At a zero, the bridge and~\eqref{eq:channel-zero} recover
a parent solution.  From a parent solution, choose the forward bridge
witnesses and set \(m\) to cancel one integral signed factor.

The full inherited interpolation, local-root, and bridge proof is kept
in \code{upstream/wave7_sun_popcount.tex}, with its precise primary
references~\cite{sun2021,sun2026}.  The root verifier checks the norm
identity by exact polynomial multiplication; the unbounded bridge is
an identified mathematical dependency.

\subsection{Exact degree and circuit cost}

The source registers have degrees
\[
 \deg X=8964,\qquad\deg Y=4713,\qquad\deg\mathcal C=217.
\]
Representative downstream degrees are
\[
\begin{array}{c|rrrrrrr}
 &A&C&P&K&D&E&F\\\hline
\deg&18393&18394&32071&32072&73574&110364&257514
\end{array}
\]
and
\[
\begin{array}{c|rrrrr}
 &G_*&H_*&I_*&\Sigma&\Theta_*\\\hline
\deg&275907&367879&1287572&91967&91980.
\end{array}
\]
The exact leading-coefficient certificate in
\code{generated/ten_integer_degree.json} propagates the degree bound
through every gate and exhibits a strictly positive final diagonal
leader.  The independent modular leading residues in
Section~\ref{sec:status} verify it with the actual fixed coefficients.
The resulting exact degree is \(8\,317\,400\).  There are 544 live
binary operations: 174 additions, 26 subtractions, and 344 multiplications.
The baseline comparison concerns unrestricted integer zeros throughout.
